% Generated draft: requires mathematical and source review.
\documentclass[11pt]{article}
\usepackage[margin=25mm]{geometry}
\usepackage{fontspec}
\setmainfont{Latin Modern Roman}
\usepackage{amsmath,amssymb,mathtools,unicode-math}
\setmathfont{Latin Modern Math}
\usepackage{enumitem,needspace,xurl,hyperref}
\setlength{\parindent}{0pt}
\setlength{\parskip}{5pt}
\setlength{\emergencystretch}{3em}
\newcommand{\N}{\mathbb{N}}
\newcommand{\Z}{\mathbb{Z}}
\newcommand{\Q}{\mathbb{Q}}
\newcommand{\R}{\mathbb{R}}
\newcommand{\C}{\mathbb{C}}
\newcommand{\F}{\mathbb{F}}
\newcommand{\A}{\mathbb{A}}
\newcommand{\PP}{\mathbb{P}}
\newcommand{\E}{\mathbb{E}}
\newcommand{\eps}{\varepsilon}
\newcommand{\dd}{\,\mathrm{d}}
\begin{document}
\section*{UnsolvedMath ID: 30003964}
\textbf{Source problem code:} catalogue entry
\subsection{Definitions and mathematical statement}
A \emph{square-tiled surface} is a compact connected translation surface \((X,\omega)\) obtained by gluing finitely many unit Euclidean squares along parallel edges by translations. Equivalently, it admits a finite translation covering of \(\C/(\Z+i\Z)\), branched over at most the image of \(0\). Here \(\omega\) is a nonzero holomorphic one-form on \(X\). Let \(\mathcal S_3\) denote the collection of such surfaces with \(g(X)=3\); an equivalence relation and a counting convention for this collection have not been specified by the supplied record.

A cylinder in a direction is a maximal open Euclidean cylinder foliated by parallel closed regular geodesics in that direction. A \emph{one-cylinder decomposition} means that a single such cylinder fills the surface apart from its boundary saddle connections. Under the interpretation that the exclusion concerns every direction, put
\[
\mathcal S_3^{\mathrm{no\,1c}}
=\{(X,\omega)\in\mathcal S_3:\text{no direction admits a one-cylinder decomposition}\}.
\]
The record does not explicitly distinguish this interpretation from exclusion in a designated direction, so the directional quantifier requires confirmation.

Kontsevich--Zorich monodromy is the homological monodromy associated with the Gauss--Manin local system over the relevant orbit space; for a square-tiled surface, one considers its associated Teichm\"uller curve, with orbifold monodromy or a suitable finite cover. Write \(\Gamma_X\leq\operatorname{GL}(L_X)\) for the image on a specified invariant integral homology lattice \(L_X\). The supplied record does not specify this representation. In particular, full homology \(H_1(X,\Z)\) has rank \(6\) in genus \(3\), whereas the reduced lattice
\[
H_1^{(0)}(X,\Z)
=\ker\bigl(p_*:H_1(X,\Z)\longrightarrow H_1(\C/(\Z+i\Z),\Z)\bigr)
\]
for a square-tiled covering \(p\) has rank \(4\). These give distinct possible interpretations of the monodromy in the question; neither is selected here.

For a specified algebraic group \(G\leq\operatorname{GL}(L\otimes_{\Z}\Q)\) defined over \(\Q\), a subgroup \(\Gamma\leq G(\Q)\cap\operatorname{GL}(L)\) is called \emph{thin relative to \(G\) and \(L\)} if it is Zariski dense in \(G\) and
\[
\bigl[G(\Q)\cap\operatorname{GL}(L):\Gamma\bigr]=\infty.
\]
One convention takes \(G\) to be the Zariski closure of \(\Gamma\); another possible intended target is a specified ambient symplectic group. The record supplies no ambient group or convention.

\textbf{Problem, with unresolved choices displayed.} Determine how frequently Kontsevich--Zorich monodromy groups are thin, with particular attention to genus-three square-tiled surfaces. In particular, after fixing the homological representation, lattice, ambient algebraic group, and the directional interpretation above, determine whether
\[
\exists (X,\omega)\in\mathcal S_3^{\mathrm{no\,1c}}
\quad\text{such that}\quad
\Gamma_X\text{ is thin relative to }G_X\text{ and }L_X.
\]
This formula states the existence target under the every-direction interpretation. It is a conditional formulation, since the missing choices have not been resolved. The broader frequency question also requires a specified population, equivalence relation, size bound, and counting measure or probability measure; no particular density or asymptotic assertion is determined by the source record.
\subsection{Short English statement}
The problem asks how common thin Kontsevich--Zorich monodromy is, especially for square-tiled surfaces of genus \(3\), and whether an example can occur when there is no one-cylinder decomposition. A precise existence statement requires specifying the homology representation, the ambient group used to define thinness, and whether the cylinder exclusion applies in every direction. A precise frequency statement additionally requires a counting or measure convention. These choices remain unresolved in the supplied record.
\subsection{Sources}
\begin{itemize}
\item Supplied local UnsolvedMath registry record, permanent ID \(30003964\), title \emph{Thin Kontsevich--Zorich Monodromy Groups}, in the Oberwolfach Reports Open Problems collection; exact source problem code \texttt{catalogue entry}. The record gives proposed year \(2018\) and supplies no named proposer.
\item Supplied report citation: \emph{Flat Surfaces and Algebraic Curves (2019)}, Oberwolfach Reports, pp.\ 2583--2650, DOI \texttt{10.4171/owr/2018/42}. Supplied source link and literature-source label ``Oberwolfach Report 42, thin monodromy'': \url{https://doi.org/10.4171/owr/2018/42}. The citation's parenthetical year \(2019\), the proposed year \(2018\), and the DOI's \(2018\) component are retained as supplied; their bibliographic relationship has not been verified.
\item Supplied literature source, labeled ``Hubert--Matheus, genus-three arithmetic monodromy'': \url{https://arxiv.org/abs/1806.09855}.
\item Supplied literature source, labeled ``Bonnafoux et al., arithmetic KZ monodromies'': \url{https://arxiv.org/abs/2206.06595}.
\item The supplied literature assessment is dated exactly \texttt{2026-08-21}. It labels the problem open and reports that genus-three arithmetic examples and families are known, while the frequency and existence of thin Kontsevich--Zorich monodromies in the general genus-three problem remain unresolved. This is supplied dated curation derived from the OWR record, not an independently verified literature assessment. The supplied background expressly calls for independent verification of time-sensitive status and cited claims.
\item The registry labels the statement \texttt{corrected\_verified} and reports a manual semantic, typographic, and LaTeX review asserting that its clean statement is self-contained and preserves the target. These are supplied editorial claims; they do not resolve the formulation choices identified here or establish mathematical correctness. No external source was visited for this draft.
\end{itemize}
\end{document}
